% App F: Poincaré gauge theory — field content and conventions.
% Forward-reference target from §2 for theory-background topics
% (vierbein/spin connection, Palatini, irreducible torsion
% decomposition formalism). The sector-by-sector enumeration of
% invariants and the cross-author coupling-constant translation
% table live in App H (pgt_enumeration.tex).

\section{Poincar\'{e} gauge theory: field content and conventions}
\label{PGTBackground}

This appendix collects the field-variable formalism that §2
invokes by reference: the tetrad and spin-connection presentation
of Poincar\'{e} gauge theory in which curvature and torsion
appear as the field strengths of local Lorentz rotations and
local translations, the Palatini equivalence between the
spin-connection formulation and the affine-connection formulation
adopted in the main text, and the full irreducible decomposition
of the torsion tensor with the explicit transformation between
its irreducible squared-norms and the direct basis of scalar
contractions used computationally.

\paragraph*{Gauging the Poincar\'{e} group}
\label{FieldContent}
We follow the pedagogical line of
Blagojevi\'{c}~\cite{blagojevic2003lectures} in deriving Poincar\'{e}
gauge theory by promoting the rigid Poincar\'{e} symmetry of
Minkowski spacetime to a local one. A matter Lagrangian
$\mathcal{L}_M(\phi,\partial_k\phi)$ that is invariant under
constant Poincar\'{e} transformations fails to remain invariant
once the parameters become spacetime-dependent functions;
restoring invariance requires two compensating fields, one for
each generator type. The Lorentz generators are compensated by a
\emph{spin connection} $A^{ij}{}_\mu$, which enters the
Lorentz-covariant derivative $\nabla_\mu\phi = (\partial_\mu +
\tfrac{1}{2}A^{ij}{}_\mu\Sigma_{ij})\phi$ acting on a matter
field $\phi$ in the representation with spin matrix $\Sigma_{ij}$.
The translation generators are compensated by the
\emph{tetrad} $b^i{}_\mu$, which relates the Riemann--Cartan
metric to the Minkowski metric of the local Lorentz
frame~\cite{kibble1961lorentz,utiyama1956invariant,sciama1962physical,barker2024poincare} via
\begin{align}
  \g{_{\mu\nu}} &= \eta_{ij}\,b^i{}_\mu\,b^j{}_\nu.
  \label{TetradMetric}
\end{align}
The commutator of two covariant derivatives produces both gauge
field strengths simultaneously: the Lorentz field strength
$F^{ij}{}_{\mu\nu} = 2\partial_{[\mu}A^{ij}{}_{\nu]} +
2A^i{}_{s[\mu}A^{sj}{}_{\nu]}$ is the curvature, and the
translation field strength $F^i{}_{\mu\nu} = 2\nabla_{[\mu}b^i{}_{\nu]}$
is the torsion. Spacetime thereby acquires the
Riemann--Cartan structure $U_4$ in which curvature and torsion are
the two independent geometric obstructions: torsion is the
antisymmetric part of the affine connection, and curvature
measures the failure of parallel transport along closed
loops~\cite{blagojevic2003lectures,hehl1976spin}. The torsion-free
limit recovers the Riemannian $V_4$ of general relativity, while
the flat-curvature limit recovers Weitzenb\"{o}ck's teleparallel
geometry $T_4$; the present work inhabits the full $U_4$.

\paragraph*{Palatini formulation and the ring connection
convention}
The tetrad / spin-connection formulation just described is
equivalent to the Palatini-style formulation used in the main
text, in which the metric $\g{_{\mu\nu}}$ and the affine
connection $\MAGA{^\sigma_\mu_\nu}$ are treated as independent
variational fields~\cite{palatini1919deduzione,ferraris1982variational,barker2024palatini}.
The bridge is the tetrad
postulate~\cite{blagojevic2003lectures}, $\nabla_\mu^{(\omega+\Gamma)}
b^i{}_\nu = 0$, which fixes a one-to-one map between the spin
connection $A^{ij}{}_\mu$ and the affine connection
$\MAGA{^\sigma_\mu_\nu}$ once the tetrad is chosen. Imposing the
metricity condition $\rcD{_\mu}\g{_{\nu\lambda}} = 0$ then
forces the metric-compatible
decomposition~\cite{blagojevic2003lectures,hehl1976spin,hehlobukhov2007cartan},
\begin{align}
  \MAGA{^\sigma_\mu_\nu} &= \rCon{^\sigma_\mu_\nu}
    + K^\sigma{}_{\mu\nu},
  \label{ContortionSplit}
\end{align}
between the Levi--Civita part and the contortion --- the same
content as the spin-connection split $A^{ij}{}_\mu = \Delta^{ij}{}_\mu(b)
+ K^{ij}{}_\mu$ that follows from the tetrad postulate, restated
here in coordinate components and consistent with §2.2's
\cref{PalatiniDecomposition}. Throughout the report we follow the
Barker-lab convention~\cite{barker2024poincare} of decorating the
Levi--Civita connection with a ring, $\rCon{^\sigma_\mu_\nu}$,
to distinguish it visually from the un-decorated Riemann--Cartan
$\MAGA{^\sigma_\mu_\nu}$.

\paragraph*{Irreducible torsion decomposition}
\label{IrreducibleTorsionFull}
The torsion tensor $\MAGT{_{\alpha\beta\mu}}$ admits a unique
decomposition into three irreducible representations of the local
Lorentz group~\cite{hehl1976spin,hehl1995metric,shapiro2002torsion,obukhov2024electrodynamics}:
the trace vector $T_\mu \equiv \MAGT{^\lambda_\lambda_\mu}$ (four
components), the axial pseudovector $S^\mu \equiv \varepsilon^{\alpha\beta\gamma\mu}
\MAGT{_{\alpha\beta\gamma}}$ (four components), and a totally
traceless tensor part $q_{\alpha\beta\mu}$ (sixteen components)
that satisfies $q^\lambda{}_{\lambda\mu} = 0$ and
$\varepsilon^{\alpha\beta\gamma\delta}q_{\alpha\beta\gamma} = 0$.
The reconstruction formula
\begin{align}
  \MAGT{_{\alpha\beta\mu}}
  &= \tfrac{1}{3}\bigl(T_\beta\,\g{_{\alpha\mu}}
                 - T_\mu\,\g{_{\alpha\beta}}\bigr)
   - \tfrac{1}{6}\,\varepsilon_{\alpha\beta\mu\nu}\,S^\nu
   + q_{\alpha\beta\mu}
  \label{IrreducibleTorsionFullDecomp}
\end{align}
expresses the same content as §2.1's
\cref{IrreducibleTorsion} restated here at appendix level so
that the orthogonality relations and the basis transformation
that follow are anchored to a single equation. The three
irreducible parts are mutually orthogonal under the natural
quadratic-norm pairing.

The torsion-squared sector of \cref{QuadraticAction} can be
written either in the irreducible basis as $a\,q^2 + b\,S^2 + c\,T^2$,
where $q^2 \equiv q_{\alpha\beta\mu}q^{\alpha\beta\mu}$ and
similarly for $S^2,T^2$, or in the direct basis $\beta_1 I_1 +
\beta_2 I_2 + \beta_3 I_3$ of scalar contractions defined in
\cref{TorsionDirectBasis}. Substituting
\cref{IrreducibleTorsionFullDecomp} into the direct-basis
contractions and using the orthogonality of the irreducible parts
gives the linear map
\begin{align}
  \begin{pmatrix} a \\ b \\ c \end{pmatrix}
  &= \begin{pmatrix}
       1 & 1 & 0 \\
       -\tfrac{1}{6} & \tfrac{1}{6} & 0 \\
       \tfrac{2}{3} & \tfrac{1}{3} & 1
     \end{pmatrix}
     \begin{pmatrix} \beta_1 \\ \beta_2 \\ \beta_3 \end{pmatrix},
  \label{DirectIrreducibleTransform}
\end{align}
which is invertible and identifies $I_3 = T^2$ trivially through
$T_\mu T^\mu = \MAGT{^\alpha_\alpha_\mu}\MAGT{^\beta_\beta^\mu}$.
The direct basis is convenient for symbolic computation because
it carries no $\varepsilon$-tensor or projector contractions; the
irreducible basis is the natural one for assigning physical
degrees of freedom to spin-parity sectors. \Cref{DirectIrreducibleTransform}
is the bridge between the two pictures used throughout §2.
